% GENERATED FILE. Edit canonical lessons, then run npm run generate:books.
% canonical-source-hash: fnv1a64:eaa5ef52c62099d7
% canonical-lessons: ZH-C198-guanji, ZH-W198-guan-shut, ZH-C198-kaiji, ZH-C198-liuyan, ZH-C198-zhanxian, ZH-C198-xinhao

\chapter{Phones and Messages}
\label{ch:zh-phones-and-messages}
\hlchaptermodality{198}

\begin{chapteropening}
\textbf{By the end of this chapter:} \emph{I can switch a phone off and on, leave a message, and say the line is engaged or there is no signal.}

Complete the last lesson of chapter 198: I can switch a phone off and on, leave a message, and say the line is engaged or there is no signal.
\end{chapteropening}

\section[\texorpdfstring{\zh{关机} (guānjī)}{关机 (guānjī)}]{\zh{关机} (guānjī) — to switch off (a phone or a computer)}

\label{lesson:ZH-C198-guanji}



\begin{quote}
Before the new one: say the Mandarin for straight on, then the Mandarin for route.
\end{quote}

\subsection*{You'll want to know: \zh{关机}}
\textbf{\zh{关机}} — \emph{guānjī} — \textquotedblleft{}to switch off (a phone or a computer)\textquotedblright{}.

\begin{quote}
\textbf{\zh{我的手机关机了}} \emph{wǒ de shǒujī guānjī le} — My phone is switched off.
\end{quote}

The new character, in pieces: \zh{关} — six strokes: two short strokes leaning toward each other, then \zh{天} underneath. The next lesson puts it in your hand.

\begin{cousinweb}
\textbf{\zh{关机}} (\emph{guānjī}) is built from \zh{关} 'to shut' and \zh{机} 'machine', both inherited from Middle Chinese, \emph{kwaen} and \emph{kj+j} (Guangyun, in Baxter's transcription). The dictionary glosses the first as the bar that shuts a door, and explains the second as the catch that sets off a crossbow.
\end{cousinweb}

\begin{grammarlens}[title={What you've built}]
One more word for reading notices and messages.
\end{grammarlens}

\subsection*{Your turn}
\begin{itemize}
\raggedright
  \item \emph{Say it:} \emph{guānjī}
  \item \emph{Say it:} \emph{guānjī}, once more
  \item \emph{Say it:} \emph{wǒ de shǒujī guānjī le} — My phone is switched off.
  \item \emph{Recall:} say the Mandarin for straight on, then the Mandarin for route, then say \emph{guānjī} again
\end{itemize}

\begin{tcolorbox}[breakable,colback=teal!4,colframe=teal!35!black,title={Before you move on}]
What does \zh{关机} mean? (\textquotedblleft{}To switch off (a phone or a computer)\textquotedblright{}.) Say it once more.
\end{tcolorbox}
\section[\texorpdfstring{\zh{关} (guān)}{关 (guān)}]{\zh{关} — two short strokes over \zh{天}}

\label{lesson:ZH-W198-guan-shut}



\begin{quote}
The word first, then the character inside it.

\begin{itemize}
\raggedright
  \item \emph{Recall:} say \emph{guānjī} — \textbf{R1}, one lesson back, before it decays
\end{itemize}

\begin{quote}
\zh{开} \zh{天}
\end{quote}

\emph{Read it:} each one aloud, with nothing said first — \emph{kāi}, \emph{tiān}
\end{quote}

\subsection*{Script — two short strokes over \zh{天}}
\hlblockfigure{\detokenize{figures/ZH-W198-guan-shut-filmstrip.pdf}}{How \zh{关} is written, stroke by stroke}

\begin{quote}
\zh{关}
\end{quote}

\textbf{Two pieces, top before bottom.}

\textbf{The top piece, two strokes.} A dot on the left, falling to the lower right, then a short stroke on the right, falling to the lower left: the two lean toward each other.

\textbf{The bottom piece, \zh{天}, four strokes:} a bar, a longer bar, a long stroke falling to the left, and a long stroke falling to the right.

2 + 4 = six.

\subsection*{Your turn}
\begin{itemize}
\raggedright
  \item \emph{Write it:} \zh{关}, all six, with the model in front of you
  \item \emph{Cover:} write it again from memory, and count the pieces, not the strokes
\end{itemize}

\begin{tcolorbox}[breakable,colback=teal!4,colframe=teal!35!black,title={Before you move on}]
How many strokes in \zh{关}? (\textbf{Six}.) Which character is the bottom piece? (\textbf{\zh{天}}.) Which way do the two top strokes lean? (\textbf{Toward each other}.)

Source: \href{https://raw.githubusercontent.com/chanind/hanzi-writer-data/68d10a4b21150cae5e1ebbd223eed289cf32d90c/data/\%E5\%85\%B3.json}{Hanzi Writer Data \zh{关}, pinned stroke paths and medians}.
\end{tcolorbox}
\section[\texorpdfstring{\zh{开机} (kāijī)}{开机 (kāijī)}]{\zh{开机} (kāijī) — to turn a phone on; to switch on a phone or a computer}

\label{lesson:ZH-C198-kaiji}



\begin{quote}
Before the new one: say the Mandarin for route, then the Mandarin for to switch off.
\end{quote}

\subsection*{You'll want to know: \zh{开机}}
\textbf{\zh{开机}} — \emph{kāijī} — \textquotedblleft{}to turn a phone on; to switch on a phone or a computer\textquotedblright{}.

\begin{quote}
\textbf{\zh{你开机了吗}} \emph{nǐ kāijī le ma} — Have you switched your phone on?
\end{quote}

\begin{cousinweb}
\textbf{\zh{开机}} (\emph{kāijī}) is built from \zh{开} 'to open' and \zh{机} 'machine', both inherited from Middle Chinese, \emph{khoj} and \emph{kj+j} (Guangyun, in Baxter's transcription).
\end{cousinweb}

\begin{grammarlens}[title={What you've built}]
One more word for reading notices and messages.
\end{grammarlens}

\subsection*{Your turn}
\begin{itemize}
\raggedright
  \item \emph{Say it:} \emph{kāijī}
  \item \emph{Say it:} \emph{kāijī}, once more
  \item \emph{Say it:} \emph{nǐ kāijī le ma} — Have you switched your phone on?
  \item \emph{Recall:} say the Mandarin for route, then the Mandarin for to switch off, then say \emph{kāijī} again
\end{itemize}

\begin{tcolorbox}[breakable,colback=teal!4,colframe=teal!35!black,title={Before you move on}]
What does \zh{开机} mean? (\textquotedblleft{}To turn a phone on; to switch on a phone or a computer\textquotedblright{}.) Say it once more.
\end{tcolorbox}
\section[liúyán]{liúyán — to leave a message}

\label{lesson:ZH-C198-liuyan}



\begin{quote}
Before the new one: say the Mandarin for to switch off, then the Mandarin for to turn a phone on.
\end{quote}

\subsection*{You'll want to know: liúyán}
\emph{liúyán} — \textquotedblleft{}to leave a message\textquotedblright{}.

\begin{quote}
\emph{tā bú zài nǐ kěyǐ liúyán} — He isn't in, but you can leave a message.
\end{quote}

Its characters come later; for now it is a word for the ear.

\begin{cousinweb}
\emph{liúyán} is built from \emph{liú} 'to leave behind' and \emph{yán} 'words', both inherited from Middle Chinese, \emph{ljuw} and \emph{ngjon} (Guangyun, in Baxter's transcription), which glosses the first 'to stay' and 'to stop', and the second 'speech'.
\end{cousinweb}

\begin{grammarlens}[title={What you've built}]
One more word for reading notices and messages.
\end{grammarlens}

\subsection*{Your turn}
\begin{itemize}
\raggedright
  \item \emph{Say it:} \emph{liúyán}
  \item \emph{Say it:} \emph{liúyán}, once more
  \item \emph{Say it:} \emph{tā bú zài nǐ kěyǐ liúyán} — He isn't in, but you can leave a message.
  \item \emph{Recall:} say the Mandarin for to switch off, then the Mandarin for to turn a phone on, then say \emph{liúyán} again
\end{itemize}

\begin{tcolorbox}[breakable,colback=teal!4,colframe=teal!35!black,title={Before you move on}]
What does liúyán mean? (\textquotedblleft{}To leave a message\textquotedblright{}.) Say it once more.
\end{tcolorbox}
\section[zhànxiàn]{zhànxiàn — engaged; busy (of a phone line)}

\label{lesson:ZH-C198-zhanxian}



\begin{quote}
Before the new one: say the Mandarin for to turn a phone on, then the Mandarin for to leave a message.
\end{quote}

\subsection*{You'll want to know: zhànxiàn}
\emph{zhànxiàn} — \textquotedblleft{}engaged; busy (of a phone line)\textquotedblright{}.

\begin{quote}
\emph{tā de diànhuà yìzhí zhànxiàn} — His line has been engaged the whole time.
\end{quote}

Its characters come later; for now it is a word for the ear.

\begin{cousinweb}
\emph{zhànxiàn} is built from \emph{zhàn} 'to take up' and \emph{xiàn} 'a line'. The second is inherited from Middle Chinese \emph{sjenH} (Guangyun, in Baxter's transcription), glossed 'thread'; the dictionary lists the first only with another sense, so how it came to mean taking up is not settled.
\end{cousinweb}

\begin{grammarlens}[title={What you've built}]
One more word for reading notices and messages.
\end{grammarlens}

\subsection*{Your turn}
\begin{itemize}
\raggedright
  \item \emph{Say it:} \emph{zhànxiàn}
  \item \emph{Say it:} \emph{zhànxiàn}, once more
  \item \emph{Say it:} \emph{tā de diànhuà yìzhí zhànxiàn} — His line has been engaged the whole time.
  \item \emph{Recall:} say the Mandarin for to turn a phone on, then the Mandarin for to leave a message, then say \emph{zhànxiàn} again
\end{itemize}

\begin{tcolorbox}[breakable,colback=teal!4,colframe=teal!35!black,title={Before you move on}]
What does zhànxiàn mean? (\textquotedblleft{}Engaged; busy (of a phone line)\textquotedblright{}.) Say it once more.
\end{tcolorbox}
\section[xìnhào]{xìnhào — phone signal}

\label{lesson:ZH-C198-xinhao}



\begin{quote}
Before the new one: say the Mandarin for to leave a message, then the Mandarin for engaged (of a phone line).
\end{quote}

\subsection*{You'll want to know: xìnhào}
\emph{xìnhào} — \textquotedblleft{}phone signal\textquotedblright{}.

\begin{quote}
\emph{zhèr méiyǒu xìnhào} — There's no signal here.
\end{quote}

Its characters come later; for now it is a word for the ear.

\begin{cousinweb}
\emph{xìnhào} is built from \emph{xìn} 'a sign, a message' and \zh{号} 'a call, a mark', both inherited from Middle Chinese, \emph{sinH} and \emph{hawH} (Guangyun, in Baxter's transcription). The dictionary glosses the first 'true' and 'proof', and the second 'to call'.
\end{cousinweb}

\begin{grammarlens}[title={What you've built}]
One more word for reading notices and messages.
\end{grammarlens}

\subsection*{Your turn}
\begin{itemize}
\raggedright
  \item \emph{Say it:} \emph{xìnhào}
  \item \emph{Say it:} \emph{xìnhào}, once more
  \item \emph{Say it:} \emph{zhèr méiyǒu xìnhào} — There's no signal here.
  \item \emph{Recall:} say the Mandarin for to leave a message, then the Mandarin for engaged (of a phone line), then say \emph{xìnhào} again
\end{itemize}

\begin{tcolorbox}[breakable,colback=teal!4,colframe=teal!35!black,title={Before you move on}]
What does xìnhào mean? (\textquotedblleft{}Phone signal\textquotedblright{}.) Say it once more.
\end{tcolorbox}
